% GENERATED FILE. Edit canonical lessons, then run npm run generate:books.
% canonical-source-hash: fnv1a64:750474e5f3c57e4e
% canonical-lessons: KA-C122-palya, KA-C122-uppinakayi, KA-C122-motte, KA-C122-minu, KA-C122-mamsa

\chapter{Dry vegetable dish, Pickle, Egg, Fish, Meat}
\label{ch:ka-dry-vegetable-dish-pickle-egg-fish-meat}
\hlchaptermodality{122}

\begin{chapteropening}
\textbf{By the end of this chapter:} \emph{I can say a dry vegetable dish, a pickle, an egg, a fish and meat.}

Complete the last lesson of chapter 122: I can say a dry vegetable dish, a pickle, an egg, a fish and meat.
\end{chapteropening}

\section[\texorpdfstring{\kn{ಪಲ್ಯ} (palya)}{ಪಲ್ಯ (palya)}]{\kn{ಪಲ್ಯ} (palya) — a dry vegetable dish}

\label{lesson:KA-C122-palya}



\begin{quote}
Before the new one: say the Kannada for rasam, then the Kannada for sambar.
\end{quote}

\subsection*{You'll want to know: \kn{ಪಲ್ಯ}}
\textbf{\kn{ಪಲ್ಯ}} — \emph{palya} — \textquotedblleft{}a dry vegetable dish\textquotedblright{}.

\begin{grammarlens}[title={What you've built}]
One more word for this chapter.
\end{grammarlens}

\subsection*{Your turn}
\begin{itemize}
\raggedright
  \item \emph{Say it:} \emph{palya}
  \item \emph{Say it:} \emph{palya}, once more
  \item \emph{Say it:} say \emph{palya}
  \item \emph{Recall:} say the Kannada for rasam, then the Kannada for sambar, then say \emph{palya} again
\end{itemize}

\begin{tcolorbox}[breakable,colback=teal!4,colframe=teal!35!black,title={Before you move on}]
What does \kn{ಪಲ್ಯ} mean? (\textquotedblleft{}A dry vegetable dish\textquotedblright{}.) Say it once more.
\end{tcolorbox}
\section[\texorpdfstring{\kn{ಉಪ್ಪಿನಕಾಯಿ} (uppinakāyi)}{ಉಪ್ಪಿನಕಾಯಿ (uppinakāyi)}]{\kn{ಉಪ್ಪಿನಕಾಯಿ} (uppinakāyi) — a pickle}

\label{lesson:KA-C122-uppinakayi}



\begin{quote}
Before the new one: say the Kannada for sambar, then the Kannada for a dry vegetable dish.
\end{quote}

\subsection*{You'll want to know: \kn{ಉಪ್ಪಿನಕಾಯಿ}}
\textbf{\kn{ಉಪ್ಪಿನಕಾಯಿ}} — \emph{uppinakāyi} — \textquotedblleft{}a pickle\textquotedblright{}.

\begin{grammarlens}[title={What you've built}]
One more word for this chapter.
\end{grammarlens}

\subsection*{Your turn}
\begin{itemize}
\raggedright
  \item \emph{Say it:} \emph{uppinakāyi}
  \item \emph{Say it:} \emph{uppinakāyi}, once more
  \item \emph{Say it:} say \emph{uppinakāyi}
  \item \emph{Recall:} say the Kannada for sambar, then the Kannada for a dry vegetable dish, then say \emph{uppinakāyi} again
\end{itemize}

\begin{tcolorbox}[breakable,colback=teal!4,colframe=teal!35!black,title={Before you move on}]
What does \kn{ಉಪ್ಪಿನಕಾಯಿ} mean? (\textquotedblleft{}A pickle\textquotedblright{}.) Say it once more.
\end{tcolorbox}
\section[\texorpdfstring{\kn{ಮೊಟ್ಟೆ} (moṭṭe)}{ಮೊಟ್ಟೆ (moṭṭe)}]{\kn{ಮೊಟ್ಟೆ} (moṭṭe) — an egg}

\label{lesson:KA-C122-motte}



\begin{quote}
Before the new one: say the Kannada for a dry vegetable dish, then the Kannada for a pickle.
\end{quote}

\subsection*{You'll want to know: \kn{ಮೊಟ್ಟೆ}}
\textbf{\kn{ಮೊಟ್ಟೆ}} — \emph{moṭṭe} — \textquotedblleft{}an egg\textquotedblright{}.

\begin{grammarlens}[title={What you've built}]
One more word for this chapter.
\end{grammarlens}

\subsection*{Your turn}
\begin{itemize}
\raggedright
  \item \emph{Say it:} \emph{moṭṭe}
  \item \emph{Say it:} \emph{moṭṭe}, once more
  \item \emph{Say it:} say \emph{moṭṭe}
  \item \emph{Recall:} say the Kannada for a dry vegetable dish, then the Kannada for a pickle, then say \emph{moṭṭe} again
\end{itemize}

\begin{tcolorbox}[breakable,colback=teal!4,colframe=teal!35!black,title={Before you move on}]
What does \kn{ಮೊಟ್ಟೆ} mean? (\textquotedblleft{}An egg\textquotedblright{}.) Say it once more.
\end{tcolorbox}
\section[\texorpdfstring{\kn{ಮೀನು} (mīnu)}{ಮೀನು (mīnu)}]{\kn{ಮೀನು} (mīnu) — a fish}

\label{lesson:KA-C122-minu}



\begin{quote}
Before the new one: say the Kannada for a pickle, then the Kannada for an egg.
\end{quote}

\subsection*{You'll want to know: \kn{ಮೀನು}}
\textbf{\kn{ಮೀನು}} — \emph{mīnu} — \textquotedblleft{}a fish\textquotedblright{}.

\begin{grammarlens}[title={What you've built}]
One more word for this chapter.
\end{grammarlens}

\subsection*{Your turn}
\begin{itemize}
\raggedright
  \item \emph{Say it:} \emph{mīnu}
  \item \emph{Say it:} \emph{mīnu}, once more
  \item \emph{Say it:} say \emph{mīnu}
  \item \emph{Recall:} say the Kannada for a pickle, then the Kannada for an egg, then say \emph{mīnu} again
\end{itemize}

\begin{tcolorbox}[breakable,colback=teal!4,colframe=teal!35!black,title={Before you move on}]
What does \kn{ಮೀನು} mean? (\textquotedblleft{}A fish\textquotedblright{}.) Say it once more.
\end{tcolorbox}
\section[\texorpdfstring{\kn{ಮಾಂಸ} (māṁsa)}{ಮಾಂಸ (māṁsa)}]{\kn{ಮಾಂಸ} (māṁsa) — meat}

\label{lesson:KA-C122-mamsa}



\begin{quote}
Before the new one: say the Kannada for an egg, then the Kannada for a fish.
\end{quote}

\subsection*{You'll want to know: \kn{ಮಾಂಸ}}
\textbf{\kn{ಮಾಂಸ}} — \emph{māṁsa} — \textquotedblleft{}meat\textquotedblright{}.

\begin{grammarlens}[title={What you've built}]
One more word for this chapter.
\end{grammarlens}

\subsection*{Your turn}
\begin{itemize}
\raggedright
  \item \emph{Say it:} \emph{māṁsa}
  \item \emph{Say it:} \emph{māṁsa}, once more
  \item \emph{Say it:} say \emph{māṁsa}
  \item \emph{Recall:} say the Kannada for an egg, then the Kannada for a fish, then say \emph{māṁsa} again
\end{itemize}

\begin{tcolorbox}[breakable,colback=teal!4,colframe=teal!35!black,title={Before you move on}]
What does \kn{ಮಾಂಸ} mean? (\textquotedblleft{}Meat\textquotedblright{}.) Say it once more.
\end{tcolorbox}
